\section{SpatialSceneQA}

Training a 3D AV-LLM demands large-scale supervision with tightly synchronized (i) metric geometry (RGB-D with camera calibration), (ii) multi-channel spatial audio, and (iii) 3D object-level annotations in a consistent reference frame. 
To this end, we introduce \textsc{SpatialSceneQA} 61K, a synthetically generated dataset for high-fidelity 3D audio-visual instruction tuning. 
Each sample contains synchronized RGB-D renderings, 4-channel FOA, and precise 3D annotations, enabling supervision for both spatial perception and audio-visual spatial reasoning.

\begin{table*}[t]
    \centering
    \caption{Statistics of the \textsc{SpatialSceneQA} 61K dataset. Tasks are grouped into ``Perception'' and ``Reasoning categories. \textbf{\#Sources} denotes active audio sources for A--B and inserted loudspeaker assets for C--E. Regarding definitions of different types of tasks, \textbf{A--B} estimate source azimuth and elevation; \textbf{C} predicts the precise Bounding Box of the inserted speaker; \textbf{D--E} identify the correspondence between a target audio speaker and a visual loudspeaker based on spatial audio-visual cues.}

    \label{tab:dataset_stats}
    
    \resizebox{\textwidth}{!}{
    \setlength{\tabcolsep}{3pt}
    \begin{tabular}{l c c l}
        \toprule
        \textbf{Task Type} & \textbf{\#Sources} & \textbf{\#Samples} & \textbf{Question \& Answer Example} \\
        \midrule
        
        % === Category 1: Perception ===
        \rowcolor{Gray} 
        \multicolumn{4}{l}{\textbf{\textit{Perception}}} \\
        \addlinespace[0.5ex]
        
        \multirow{2}{*}{\textbf{A: Single-Source Audio DoA}} & \multirow{2}{*}{1} & \multirow{2}{*}{32K} & 
        \textbf{Q:} Based on the audio cues, please output the precise azimuth and elevation. \\
        & & & \textbf{A:} azimuth: -7; elevation: -22 \\
        % \addlinespace[1ex]

        \midrule
        
        \multirow{2}{*}{\textbf{B: Overlap-Source Audio DoA}} & \multirow{2}{*}{2} & \multirow{2}{*}{30K} & 
        \textbf{Q:} Where is the female voice coming from? Output azimuth and elevation. \\
        & & & \textbf{A:} azimuth: 28; elevation: -15 \\
        
        \midrule
        
        \multirow{3}{*}{\textbf{C: Visual Grounding}} 
        & 1 & 17K & \multirow{3}{*}{\shortstack[l]{
            \textbf{Q:} Identify the 3D box for the speaker in this scene. \\
            \textbf{A:} $\texttt{bbox\_0} = \texttt{Bbox}(\text{speaker}, 0.14, -0.48, -1.15, 0.33, 0.88, 0.32)$ 
        }} \\
        & 2 & 34K & \\
        & 3 & \phantom{0}9K & \\
        \midrule[\heavyrulewidth]
        % === Category 2: Reasoning ===
        \rowcolor{Gray} 
        \multicolumn{4}{l}{\textbf{\textit{Reasoning}}} \\
        \addlinespace[0.5ex]
        \multirow{2}{*}{\textbf{D: Single-Source Multi-speakers Matching}} 
        & 2 & 10K & \textbf{Q:} Determine which speaker corresponds to the audio source. \\
        & 3 & \phantom{0}4K & \textbf{A:} Left \\
        % \addlinespace[1ex]
        \midrule
        \multirow{2}{*}{\textbf{E: Overlap-Source Multi-speakers Matching}} 
        & 2 & 24K & \textbf{Q:} Can you tell which speaker the male voice originates from? \\
        & 3 & \phantom{0}5K & \textbf{A:} Center \\
        \bottomrule
    \end{tabular}
    }
\end{table*}

\subsection{Data Simulation Pipeline}

We synthesize \textsc{SpatialSceneQA} using \textit{SoundSpaces 2.0} \citep{chen22soundspaces2}, which supports continuous, geometry-aware acoustic rendering in scanned 3D environments. 
The pipeline generates synchronized FOA audio and RGB-D observations, together with calibrated camera metadata and 3D ground truth, as presented in Fig.~\ref{fig:pipeline}.


\textbf{Audio Source Selection.}
We use dry monaural speech from \textit{LibriSpeech} \citep{panayotov2015librispeech}, which contains read speech recordings with aligned transcripts and speaker metadata. To prevent data leakage, we construct the training set from train-clean-100, use dev-clean for validation, and reserve test-clean for evaluation. For tasks involving speaker identity attributes, the provided gender labels are used to specify the target speaker in the instruction.

\textbf{Acoustic Simulation.} \textit{SoundSpaces 2.0} renders multi-channel spatial audio by simulating room impulse responses (RIRs) over realistic 3D meshes with material-dependent acoustics. 
It employs bidirectional path tracing \citep{cao2016interactive} to model direct sound and higher-order effects, including reflections, transmission, diffraction, and air absorption. 
Given a monaural source signal $A^{(s)}(t)$ emitted from position $\mathbf{s}\in\mathbb{R}^3$ and a receiver at position $\mathbf{r}\in\mathbb{R}^3$ with heading $\theta\in\mathbb{R}$, the received FOA signal on channel $c$ is
\begin{equation}
    A^{(r)}_c(t) = \bigl(R_c(\cdot;\mathbf{s},\mathbf{r},\theta)\ast A^{(s)}\bigr)(t),
\end{equation}
where $\ast$ denotes convolution over time, $t \in [1, T]$ and $c \in \{0, 1, 2, 3\}$ indexes time steps and FOA channels respectively, $R_c(t, \mathbf{s}, \mathbf{r}, \theta)$ stands for the channel-specific RIR.
To reduce degenerate cases caused by fully occluded direct paths, we sample source--receiver pairs within the same room. 
Consequently, we sample a navigable receiver pose, then sample a source position within 1--4\,m, enforcing a 0.5\,m clearance from obstacles. 
Pairs were retained only when the geodesic distance was less than twice the Euclidean distance, ensuring navigability and connectivity.

\textbf{Visual Scene Construction.} We use \textit{Habitat-Sim} \citep{szot2021habitat} for rendering RGB-D observations aligned to the audio. Scenes are drawn from the semantically annotated subset of the Habitat-Matterport 3D Dataset (HM3D) \citep{ramakrishnan2021habitat}, enabling precise instance-level ground truth. Scene-based data partitioning was employed, with 130/15/36 scenes for train/val/test, respectively. The synchronized RGB, depth, semantic instance masks, and calibrated camera metadata, along with the receiver pose and the ground-truth source poses, are exported.

\textbf{Synthetic 3D Object Generation.} 
Static HM3D scans contain limited diversity of visually salient sound-emitting objects. 
To increase appearance and geometry diversity while keeping metric annotations exact, synthesized loudspeaker meshes were generated by \textit{Hunyuan3D-1.0} \cite{yang2024hunyuan3d} and further inserted. 
We generate 120 distinct floor-standing loudspeaker models and split them into 96/12/12 for train/validation/test to evaluate generalization to unseen geometries. 
A visibility constraint is enforced by requiring each target loudspeaker to occupy at least 500 pixels in the semantic frame (rendered at $1920 \times 1080$), filtering out heavily occluded or distant instances.

\subsection{Creation of QA Dataset}

We convert each simulated scene into instruction-answer pairs as shown in Table~\ref{tab:dataset_stats}, spanning two modules: \textit{Perception} and \textit{Reasoning}.

\textbf{Perception Tasks (Task A--C).}
Tasks A--B supervise acoustic localization by regressing DoA in spherical coordinates (azimuth, elevation) under both single-source task (A) and overlapping-source task (B) settings.
We express angles in degrees under a right-handed camera coordinate system with $x$ pointing right, $y$ up, and $z$ pointing backwards; azimuth is measured on the horizontal plane with $0^\circ$ facing forward, and elevation is measured from the horizontal plane.

Task C targets metric 3D visual grounding of sound-emitting objects. 
Following N3D-VLM \cite{wang2025n3d}, we represent each 3D bounding box as a structured sequence $\texttt{bbox}(c, x, y, z, s_x, s_y, s_z)$, where $c$ is the category label, $(x,y,z)$ is the box center in the camera coordinate system, and $(s_x,s_y,s_z)$ are axis-aligned box dimensions. 
To avoid learning shortcuts from domain gaps, $\{1,2,3\}$ loudspeaker instances were randomly inserted per scene and require the model to output 3D boxes for all visible loudspeakers, enforcing explicit geometric grounding.

\textbf{Reasoning Tasks (Task D--E).}
Tasks D--E evaluate cross-modal spatial correspondence: given spatial audio and an RGB-D observation containing multiple candidate loudspeakers, the model must identify which visual loudspeaker matches a target speech source. 
Task D uses a single active source, while Task E introduces overlapping sources and specifies the target by gender. 
To avoid the appearence of trivial solutions, at least two visible candidate loudspeakers in every reasoning sample are guaranteed. 
Solving these tasks requires integrating (i) directional cues from FOA, (ii) 3D localization of candidate loudspeakers, and (iii) metric alignment between acoustic DoA and visual geometry, thereby directly probing audio-visual embedding alignment and multi-step spatial reasoning.
